%%%%%%%%%%%%%%%%%%%%%%%%%%%%%%%%%%%%%%%%%%%%%%%%
% Wartungsplan CNC-Fräse BZT PFX 500H (Formular zum Ausdrucken)
% CC-BY-SA 3.0, FAU FabLab
% Früher Fraese_Wartungsplan.fods (LibreOffice), jetzt LaTeX.
%%%%%%%%%%%%%%%%%%%%%%%%%%%%%%%%%%%%%%%%%%%%%%%%
\PassOptionsToPackage{table}{xcolor}
\documentclass[a4paper,landscape,10pt]{scrartcl}
\usepackage[T1]{fontenc}
\usepackage[utf8]{inputenc}
\usepackage{lmodern}
\renewcommand{\familydefault}{\sfdefault}
\usepackage[ngerman]{babel}
\usepackage[a4paper,landscape,left=12mm,right=12mm,top=10mm,bottom=16mm,footskip=8mm]{geometry}
\usepackage{graphicx}
\usepackage{xcolor}
\usepackage{amssymb}
\usepackage{tabularx}
\usepackage{lastpage}
\usepackage{fancyhdr}
\usepackage{eso-pic}

% Versionsnummer wie in fablab-document (version.tex erzeugt make)
\InputIfFileExists{version}{}{\newcommand{\Version}{Entwurf \today}}

\definecolor{fablab}{RGB}{20,51,105}
% Intervallfarben (wie im bisherigen Plan)
\definecolor{iW}{HTML}{FFCC00}   % wöchentlich / 14-tägig
\definecolor{iM}{HTML}{83CAFF}   % monatlich
\definecolor{iQ}{HTML}{9999FF}   % vierteljährlich
\definecolor{iH}{HTML}{579D1C}   % halbjährlich
\definecolor{iJ}{HTML}{FF420E}   % jährlich
\definecolor{iT}{gray}{0.85}     % täglich (nicht eintragen)

\pagestyle{fancy}
\fancyhf{}
\renewcommand{\headrulewidth}{0pt}
\fancyfoot[L]{\scriptsize Nachdrucken: \texttt{Dateifreigabe/1\_FabLab\_Einweisungen/fraese-einweisung/Fraese\_Wartungsplan.pdf}
  \quad Quelle: \texttt{github.com/fau-fablab/fraese-einweisung}}
\fancyfoot[R]{\scriptsize Version \Version\quad Seite \thepage\ von \pageref*{LastPage}}
\setlength{\parindent}{0pt}

% ---------------------------------------------------------------------
%  Bausteine
% ---------------------------------------------------------------------
\usepackage{tikz}
\usepackage{xstring}
\newcommand{\EW}{\,\textsuperscript{\textbf{EW}}}   % Details in der Einweisung

% Raster: Breiten und Höhen
\newlength{\wG}\setlength{\wG}{8mm}     % Gruppe
\newlength{\wI}\setlength{\wI}{29mm}    % Intervall
\newlength{\wT}\setlength{\wT}{74mm}    % Tätigkeit
\newlength{\wS}                         % Termin (aus der Restbreite)
\newlength{\hZ}\setlength{\hZ}{7.4mm}   % Zeilenhöhe
\newcounter{zeile}

% Eine Aufgabe:  \Aufgabe{Farbe}{Intervall}{Tätigkeit}{Muster}
%   Muster: 12 Zeichen für die Termine 1./15. Jan (Jul) ... 1./15. Jun (Dez)
%   x = fällig (Datum eintragen), . = nicht fällig, t = täglich (nicht eintragen),
%   - = leeres Feld (z.B. für handschriftliche Aufgaben)
\newcommand{\Aufgabe}[4]{%
  \stepcounter{zeile}%
  \pgfmathsetlengthmacro{\yo}{-\value{zeile}*\hZ+\hZ}%
  \pgfmathsetlengthmacro{\yu}{-\value{zeile}*\hZ}%
  \ifx\relax#1\relax\else
    \fill[#1] (\wG,\yo) rectangle ++(\wI,-\hZ);
  \fi
  % weiße Schrift auf dunklen Intervallfarben
  \def\tc{black}\IfStrEqCase{#1}{{iH}{\def\tc{white}}{iJ}{\def\tc{white}}}[]%
  \node[anchor=west, font=\small\bfseries, text=\tc]
    at (\wG+1.2mm, \yo-0.5\hZ) {#2};
  \node[anchor=west, font=\small, text width=\wT-2.4mm]
    at (\wG+\wI+1.2mm, \yo-0.5\hZ) {#3};
  \foreach \i in {1,...,12}{%
    \StrChar{#4}{\i}[\c]%
    \pgfmathsetlengthmacro{\xs}{\wG+\wI+\wT+(\i-1)*\wS}%
    \IfStrEq{\c}{x}{\fill[#1!35] (\xs,\yo) rectangle ++(\wS,-\hZ);}{}%
    \IfStrEq{\c}{.}{\fill[black!12] (\xs,\yo) rectangle ++(\wS,-\hZ);}{}%
    \IfStrEq{\c}{t}{\fill[iT!55] (\xs,\yo) rectangle ++(\wS,-\hZ);}{}%
  }%
  \draw[black!60, line width=0.3pt] (\wG,\yu) -- (\wG+\wI+\wT+12\wS,\yu);
}
% Trennlinie zwischen Gruppen
\newcommand{\Trenner}{%
  \pgfmathsetlengthmacro{\yt}{-\value{zeile}*\hZ}%
  \draw[black, line width=1.2pt] (0,\yt) -- (\wG+\wI+\wT+12\wS,\yt);}
% Gruppenbezeichnung links, von Zeile #1 bis #2
\newcommand{\Gruppe}[3]{%
  \pgfmathsetlengthmacro{\ya}{-(#1-1)*\hZ}%
  \pgfmathsetlengthmacro{\yb}{-#2*\hZ}%
  \node[rotate=90, font=\footnotesize\bfseries, align=center]
    at (0.5\wG, 0.5*\ya+0.5*\yb) {#3};}

\begin{document}

% ------------------------------ Kopf ---------------------------------
\begin{minipage}[b]{0.72\linewidth}
  {\LARGE\bfseries Wartungsplan CNC-Fräse BZT PFX 500H}\\[2mm]
  \normalsize
  Halbjahr:\enspace $\square$ 1.\ Halbjahr (Jan--Jun)\quad $\square$ 2.\ Halbjahr (Jul--Dez)
  \qquad Jahr: 20\rule{10mm}{0.4pt}
  \qquad Plan begonnen am: \rule{25mm}{0.4pt}
\end{minipage}\hfill
\begin{minipage}[b]{0.25\linewidth}\raggedleft
  \IfFileExists{fablab-document/logo/Logo/Logo-FAU-bunt.pdf}%
    {\includegraphics[height=11mm]{fablab-document/logo/Logo/Logo-FAU-bunt.pdf}}{\textbf{FAU FabLab}}
\end{minipage}\\[1mm]
{\color{fablab}\rule{\linewidth}{0.8pt}}\\[1.5mm]
\small
In die \textbf{farbigen Felder} Datum und Kürzel eintragen, wenn die Aufgabe erledigt ist.
Graue Felder: nicht fällig. Tägliche Aufgaben werden nicht eingetragen.
\textsuperscript{\textbf{EW}}\,= Details in der Einweisung (Abschnitt „Wartung“).\\[1mm]
\begin{tabular}{@{}l@{\quad}*{6}{l@{\quad}}}
  \textbf{Intervall:} &
  \colorbox{iT}{\strut\ täglich\ } &
  \colorbox{iW}{\strut\ wöchentlich / 14-tägig\ } &
  \colorbox{iM}{\strut\ monatlich\ } &
  \colorbox{iQ}{\strut\ vierteljährlich\ } &
  \colorbox{iH}{\strut\ \color{white}halbjährlich\ } &
  \colorbox{iJ}{\strut\ \color{white}jährlich\ }
\end{tabular}

\vspace{2.5mm}
% ------------------------------ Plan ---------------------------------
\setlength{\wS}{\dimexpr(\linewidth-\wG-\wI-\wT)/12\relax}
\begin{tikzpicture}[x=1pt,y=1pt]
  % Kopfzeile des Plans: Monate und Termine
  \pgfmathsetlengthmacro{\xa}{\wG+\wI+\wT}
  \node[anchor=west, font=\small\bfseries] at (1.2mm, 9mm) {Aufgabe};
  \node[anchor=west, font=\footnotesize] at (1.2mm, 3.2mm) {Intervall und Tätigkeit};
  \node[anchor=east, font=\footnotesize] at (\xa-1.2mm, 3.2mm) {Termin (Tag des Monats):};
  \foreach \m/\n [count=\k from 0] in {Jan/Jul,Feb/Aug,Mär/Sep,Apr/Okt,Mai/Nov,Jun/Dez}{
    \node[font=\small] at (\xa+2*\k*\wS+\wS, 9mm) {\textbf{\m}\,/\,\n};
    \node[font=\footnotesize] at (\xa+2*\k*\wS+0.5\wS, 3.2mm) {1.};
    \node[font=\footnotesize] at (\xa+2*\k*\wS+1.5\wS, 3.2mm) {15.};
  }
  \draw[black, line width=1.2pt] (0,0) -- (\xa+12\wS,0);
  \setcounter{zeile}{0}
  % Reinigung und Schmierung
  \Aufgabe{iT}{täglich}{Spindel: Dichtung und HSK\EW}{tttttttttttt}
  \Aufgabe{iT}{täglich}{Faltenbalg absaugen}{tttttttttttt}
  \Aufgabe{iW}{14-tägig}{HSK-Kegel einsprühen\EW}{xxxxxxxxxxxx}
  \Aufgabe{iW}{$\le$ 14-tägig}{Maschine reinigen}{xxxxxxxxxxxx}
  \Aufgabe{iW}{14-tägig}{KSS nachfüllen, wenn nötig}{xxxxxxxxxxxx}
  \Aufgabe{iM}{monatlich}{Greifer schmieren\EW}{x.x.x.x.x.x.}
  \Aufgabe{iM}{monatlich}{Maschine gründlich reinigen}{x.x.x.x.x.x.}
  \Aufgabe{iQ}{vierteljährlich}{Scheiben gründlich reinigen}{x.....x.....}
  \Aufgabe{iH}{halbjährlich}{Führung und Spindeln schmieren\EW}{....x.......}
  \Gruppe{1}{9}{Reinigung /\\ Schmierung}
  \Trenner
  % Prüfung
  \Aufgabe{iT}{täglich}{Luftaustritt an der Spindel\EW}{tttttttttttt}
  \Aufgabe{iW}{wöchentlich}{Auf Schäden prüfen}{xxxxxxxxxxxx}
  \Aufgabe{iM}{monatlich}{Kühlwasserstand prüfen}{x.x.x.x.x.x.}
  \Aufgabe{iQ}{vierteljährlich}{Scheiben auf Risse prüfen}{x.....x.....}
  \Aufgabe{iH}{halbjährlich}{Kabel und Not-Aus prüfen\EW}{........x...}
  \Aufgabe{iH}{halbjährlich}{Weitere Prüfungen\EW}{........x...}
  \Gruppe{10}{15}{Prüfung}
  \Trenner
  % Austausch
  \Aufgabe{iJ}{jährlich}{Kühlwasser tauschen {\footnotesize(nur im Juni)}}{..........x.}
  \Gruppe{16}{16}{\scriptsize Tausch}
  \Trenner
  % Sonstiges: frei für weitere Aufgaben
  \Aufgabe{}{}{}{------------}
  \Aufgabe{}{}{}{------------}
  \Gruppe{17}{18}{Sonst.}
  % Rahmen und senkrechte Linien
  \pgfmathsetlengthmacro{\yu}{-\value{zeile}*\hZ}
  \draw[black, line width=1.2pt] (0,14mm) rectangle (\xa+12\wS,\yu);
  \draw[black, line width=0.8pt] (\wG,0) -- (\wG,\yu);
  \draw[black!60, line width=0.3pt] (\wG+\wI,0) -- (\wG+\wI,\yu);
  \draw[black, line width=1.2pt] (\xa,14mm) -- (\xa,\yu);
  \foreach \i in {1,...,11}{
    \pgfmathsetlengthmacro{\xs}{\xa+\i*\wS}
    \pgfmathparse{int(mod(\i,2))}
    \ifnum\pgfmathresult=0
      \draw[black, line width=0.8pt] (\xs,14mm) -- (\xs,\yu);
    \else
      \draw[black!60, line width=0.3pt] (\xs,6.4mm) -- (\xs,\yu);
    \fi
  }
  \draw[black!60, line width=0.3pt] (\xa,6.4mm) -- (\xa+12\wS,6.4mm);
\end{tikzpicture}

\newpage
{\Large\bfseries Unregelmäßigkeiten} \hfill \small Wartungsplan CNC-Fräse BZT PFX 500H\\[2mm]
\normalsize
\setlength{\tabcolsep}{4pt}
\renewcommand{\arraystretch}{1}
\begin{tabularx}{\linewidth}{|p{22mm}|p{80mm}|X|p{20mm}|}
  \hline
  \textbf{Datum} & \textbf{Problem} & \textbf{Beschreibung / Maßnahme} & \textbf{Kürzel} \\
  \hline
  & & & \rule[-3mm]{0pt}{10.5mm}\\ \hline
  & & & \rule[-3mm]{0pt}{10.5mm}\\ \hline
  & & & \rule[-3mm]{0pt}{10.5mm}\\ \hline
  & & & \rule[-3mm]{0pt}{10.5mm}\\ \hline
  & & & \rule[-3mm]{0pt}{10.5mm}\\ \hline
  & & & \rule[-3mm]{0pt}{10.5mm}\\ \hline
  & & & \rule[-3mm]{0pt}{10.5mm}\\ \hline
  & & & \rule[-3mm]{0pt}{10.5mm}\\ \hline
  & & & \rule[-3mm]{0pt}{10.5mm}\\ \hline
  & & & \rule[-3mm]{0pt}{10.5mm}\\ \hline
  & & & \rule[-3mm]{0pt}{10.5mm}\\ \hline
  & & & \rule[-3mm]{0pt}{10.5mm}\\ \hline
  & & & \rule[-3mm]{0pt}{10.5mm}\\ \hline
  & & & \rule[-3mm]{0pt}{10.5mm}\\ \hline
  & & & \rule[-3mm]{0pt}{10.5mm}\\ \hline
  & & & \rule[-3mm]{0pt}{10.5mm}\\ \hline
\end{tabularx}

\end{document}
